\section{WC 4/11/24}

\subsection*{Overview:}
This week I selected the final form of my project. With this chosen I began to dive into research regarding some technical aspects of my project. I also began to restructure my EPQ project files and notes in preparation for them to expand (and integrate into my software project shortly in the future)

\subsection*{Achievements, Thoughts and Ideas:}
I started exploring the more artistic aspects of my project including \nameref{research:UIDesign} and \nameref{research:ColourTheory}. While this is early I wanted to form some fundamental ideas that would act as the foundations from which the software would build upon. I want this project to use the software I build as a tangible representation to the research I achieve with regard to psychology and human learning.

I began thinking about the fundamental principles which would dictate how my software would function and how this would benefit my users. I undertook a \nameref{research:MindMaps} research exploring my initial ideas, but these ideas will be fluid and changing as I receive feedback from industry experts and peers alike (the latter being highly important as they make up my primary target demographic).

I then began making early software prototypes to discover out what I had to learn from gaps in my knowledge, no visuals to show (yet) but I built my early database and worked on backend code (the platform that the user's device will communicate with to store their data). This will likely be rewritten very shortly but taught me a lot about both what works (and more importantly) what doesn't thus improving my skills.

I began collating rolling lists of annoyances I find in daily note-taking and problems I wish to solve:\nameref{research:NotetakingAnnoyances}, this will help me find challenges that my end artifact can solve further improving the value of my project.

I then watched a webinar on finding good research which, in addition to being useful for finding future sources, highlighted some other pre-existing tools I can learn from whilst building my project.



\subsection*{What went well?}
This week brings the very first section of my EPQ to a close, with a project now chosen and goals set I can work on achieving many more smaller micro goals that will help build up a project I am proud of.

\subsection*{What did I learn?}

Whilst conducting research in many, many fields to establish technical feasibility of my design, I covered the foundations of many topics. While this hasn't resulted in many research reports, yet I am sure these learnings will prove incredibly useful in the not too distant future.


\subsection*{Plans for next week:}
Next week I plan to have completed some visual demos that I can use to try and gauge the scope of this project.  With demos in tow I can get feedback from friends and family to help improve my base ideas before I spend time building up a final Artefact.
